\section{Introduction}\label{sec:introduction}

Autonomous close-range tracking and intercept guidance for unmanned aerial vehicles (UAVs) is a critical capability in defense applications, including target escort, threat interception, and cooperative surveillance~\cite{pn}. In terminal engagement scenarios, the pursuer must maintain a prescribed aspect angle and minimum stand-off distance while the target may execute evasive maneuvers. Classical proportional navigation (PN) and its augmented variants provide analytical guarantees under idealized kinematic assumptions~\cite{pn}, but their performance degrades when target dynamics deviate from the expected model or when engagement geometries become unfavorable.

Deep reinforcement learning (DRL) offers an appealing alternative: rather than hand-designing guidance laws, one can train a neural policy to directly map sensor observations to control commands~\cite{ppo,sb3}. End-to-end DRL approaches have demonstrated impressive results in simulated aerial combat and racing tasks~\cite{mn2019,ernest2019}. However, a fundamental tension remains. End-to-end policies must simultaneously learn perception, prediction, guidance, and low-level control---a burden that demands enormous sample complexity and often produces policies that overfit to training conditions~\cite{mn2019}. In safety-critical guidance tasks, where sim-to-real transfer is essential, this sample-efficiency gap is a decisive limitation.

A hierarchical decomposition naturally suggests itself: separate the guidance problem into (1) a high-level strategy that optionally selects a virtual pursuit point (VPP) relative to the target, (2) a classical line-of-sight-rate (LOS-rate) guidance law that generates roll-rate, normal-overload, and throttle commands, and (3) a low-level controller that stabilizes the aircraft attitude~\cite{mpn2022}. This architecture reduces the policy's responsibility to high-level geometric decision-making while preserving the analytical rigor of classical guidance for command generation.

Despite the intuitive appeal of hierarchical guidance, the literature lacks rigorous evidence about \emph{which} architectural choice is the dominant driver of performance. Is the learned VPP offset essential, or does the guidance law itself provide most of the benefit? Is end-to-end DRL merely harder to tune, or fundamentally unable to match a hierarchical design? We address these questions through a controlled 5-seed fair comparison across three methods---hierarchical VPP+LOS-rate, zero-offset (No-VPP), and end-to-end (E2E) DRL---evaluated on the JSBSim F-16 high-fidelity aerodynamic backend with 1440 total episodes across four canonical engagement scenarios. Our contributions are fourfold:

\begin{enumerate}
    \item We establish the \textbf{quantitative advantage of the hierarchical VPP architecture} on the disadvantage engagement geometry. VPP is the \emph{only} method to achieve non-zero success ($20\%$ [12,28] 95\% bootstrap CI), while No-VPP crashes at $100\%$ (Fisher exact $p = 6.5 \times 10^{-37}$ for crash-rate reduction) and E2E crashes at $75\%$. On the remaining three scenarios (favorable, neutral, challenging), VPP and No-VPP both achieve $100\%$ success, while E2E reaches only $87.5\%$ on neutral ($p = 5.0 \times 10^{-5}$) and $75\%$ on challenging ($p = 4.8 \times 10^{-10}$).

    \item We provide \textbf{mechanistic evidence for VPP's learned safety-valve behavior}. The VPP policy's mean virtual-point shift on disadvantage is 2987~m---3.2$\times$ larger than on challenging (936~m) and 2.9$\times$ larger than on neutral (1048~m)---demonstrating that the policy adaptively selects large forward offsets in high-risk geometries to maintain safe stand-off distance, trading timeout (60\%, $-300$ reward) for crash avoidance ($-600$ reward). This is an emergent conservative strategy consistent with the reward structure, not an architecture limitation.

    \item We confirm that \textbf{CEM-EMA gain optimization is a major performance driver} ($+14.47$ percentage points on the simple backend), but its benefits are architecture-conditional: without the VPP safety margin, even optimally-tuned guidance gains cannot prevent $100\%$ crash on the disadvantage geometry under JSBSim aerodynamics.

    \item We demonstrate that the \textbf{hierarchical decomposition is strictly more capable than E2E DRL} under identical training budgets. After extending E2E training to 2~M steps (10$\times$ the hierarchical budget) with aligned reward functions, E2E remains unstable and fails to match the hierarchical design's success rate on neutral and challenging scenarios.
\end{enumerate}

The remainder of the paper is organized as follows. Section~\ref{sec:related_work} reviews classical guidance, data-driven methods, and hybrid approaches. Section~\ref{sec:methods} describes the hierarchical architecture, bilevel optimization framework, training protocol, and scenario design. Section~\ref{sec:results} presents the 5-seed JSBSim fair comparison, architecture ablations, and CEM optimizer analysis. Section~\ref{sec:discussion} interprets the findings, with emphasis on the VPP safety-valve mechanism and the architecture-conditional nature of gain optimization. Section~\ref{sec:limitation} discusses limitations and future directions. Section~\ref{sec:conclusion} summarizes conclusions.
